\cotefile{136}
\scanpage{330}

\diagram{\begin{tikzcd}
\Gamma^p(\Gamma^q(X)) \\
\Gamma^{pq}(X)\arrow[u,"{\deleted{\beta}\inserted{\gamma}^{X}_{p,q}}"']
\end{tikzcd}}

\diagram{\begin{tikzcd}
\beta^X_{p,q}:\ \Gamma^p(X)\otimes\Gamma^q(X)\arrow[r] & \Gamma^{pq}(X) \\
\alpha^{X,Y}_{pq}:\ \Gamma^{\unsure{*|p}}(X)\otimes\Gamma^*(Y)\arrow[r,"\sim"] & \Gamma^{2}(X\times Y)
\end{tikzcd}}

\[\Delta:\ A\longleftarrow A\otimes A\]
\[u_p:\ A\longleftarrow\Gamma^p A\]
\diagram{\begin{tikzcd}
F(X\otimes Y) & F(X)\otimes F(Y)\arrow[l,"\tau_{X,Y}"'] \\
F(\Gamma^p(X)) & \Gamma^pF(X)\arrow[l,"\sigma^{\alpha}_p"']
\end{tikzcd}}
OK\quad OK

\par\noindent\rule{.9\linewidth}{.3pt}

\deleted{Soit $A$ algèbre sur un anneau $k$ \deleted{de car. $0$}, on va définir un hom can. \deleted{$k$-lin.}
\[
\operatorname{Sym}^*A\longrightarrow A
\]
par ses restrictions
\[
\operatorname{Sym}^n(A)\xrightarrow{u_n} A\qquad(n\geq0)
\]
et en prenant
\[
u_n(\underbrace{x_1.x_2.\ldots.x_n}_{\substack{\text{produit}\\\text{dans }\operatorname{Sym}^*(A)}})
=\frac1{n!}x_1x_2\ldots x_n
\]
Je dis que le diagramme
\diagram{\begin{tikzcd}[column sep=large,row sep=large]
\operatorname{Sym}^{p+q}A\arrow[r,"C_{pq}u_{p+q}"]\arrow[d,"\Delta"'] & A \\
\operatorname{Sym}^pA\otimes\operatorname{Sym}^qA\arrow[r,"u_p\otimes u_q"'] & A\otimes A\arrow[u,"\text{mult. de }A"']
\end{tikzcd}}

i.e un hom can. de
\[
\deleted{(\Gamma^\alpha(A)\otimes A)^{\wedge}}
\]
\[
\deleted{\Gamma^\alpha,\,V(A)}\operatorname{Hom}_{k\text{-lin}}(\operatorname{Sym}^*A,B
\]
\[
=\Gamma^*(\check A)\otimes_k B\qquad\text{i.e.}\ \text{un élt}\ \text{de}
\]
\[
\Gamma^*(\check A)\otimes B\qquad\text{p. ex.}\text{ de }\check A\otimes_k B
\]
\deleted{\unsure{$u$}}
\[
A\underset{\unsure{h|k}}{\Longrightarrow} B\qquad k\text{ lin.}
\]
\[
\operatorname{Sym}^* A
\]
\[
\operatorname{Sym}^*(
\]
}

\mspage{331}
